% sections/abstract-en.tex --- English abstract + keywords (9 pt, 1 cm indent).

\begin{citabstract}
\noindent\textbf{Bohusevych O., Svystunov A. PARCS-Agent: An MCP-Mediated Bridge from LLM Agents to Distributed Compute Clusters.}
Contemporary large language model (LLM) agents invoke external tools for information retrieval, code execution, and application-programming-interface access; each such invocation runs, however, as a local synchronous function inside a single process. Compute-bound subtasks (Monte Carlo simulation, hyperparameter exploration, combinatorial search) are therefore constrained to the agent's single execution thread, and no standardised mechanism exists for an agent to acquire additional computational resources during task execution. This paper proposes PARCS-Agent, an architectural pattern in which a server implementing the Model Context Protocol (MCP) provides an interface to a Kubernetes-based PARCS distributed computing cluster and exposes a compact set of tools enabling the agent to submit arbitrary C\# computation logic at runtime. The platform compiles the submitted source using Roslyn, distributes execution across daemon worker pods through Kubernetes Event-Driven Autoscaling (KEDA), and propagates intermediate results between successive parallel layers. The agent generates only per-worker logic in straight-line code, thereby circumventing the documented limitations of LLMs in producing correct MPI, OpenMP, and CUDA programs. Bulk input data is delivered through a content-addressed shared-storage side channel, decoupling dataset volume from protocol payload. The approach is illustrated by a Monte Carlo Value-at-Risk case study; quantitative evaluation is reserved for a subsequent publication.

\medskip
\noindent\textbf{Keywords:} AI agents; large language models; Model Context Protocol;
distributed computing; parallel computation; LLM tool use; Kubernetes;
agent infrastructure.
\end{citabstract}
